% !TeX program = xelatex
\documentclass[11pt,aspectratio=43,compress]{beamer}
\usepackage{ruc-research}

% 方括号内为页脚短标题、姓名和机构。
\title[汇报标题]{汇报标题}
\subtitle{副标题（可选）}
\author[汇报人]{汇报人姓名}
\institute[机构]{机构名称}
\date{汇报日期}

\begin{document}

\begin{frame}[plain,noframenumbering]
  \titlepage
\end{frame}

% 删除目录页后，\section 仍可用于顶部导航。
\begin{frame}{目录}
  \tableofcontents
\end{frame}

\section{版式示例}

\begin{frame}{分栏与文本框}
  \begin{columns}[T,onlytextwidth]
    \begin{column}{.47\textwidth}
      \textbf{正文栏}
      \begin{itemize}
        \item 一级条目
        \item \alert{重点文字}
      \end{itemize}
    \end{column}
    \begin{column}{.47\textwidth}
      \begin{block}{文本框标题}
        在此填写需要单独展示的文字。
      \end{block}
    \end{column}
  \end{columns}
\end{frame}

\begin{frame}{公式与交叉引用}
  \begin{equation}\label{eq:example}
    y = \alpha + \beta x.
  \end{equation}
  式~\eqref{eq:example} 中，$\alpha$ 为截距，$\beta$ 为斜率。
\end{frame}

\section{图表示例}

\begin{frame}{三线表}
  \begin{table}
    \centering
    \caption{表格标题}\label{tab:example}
    \begin{tabularx}{.84\textwidth}{@{}Yrr@{}}
      \toprule
      项目 & 组 A & 组 B \\
      \midrule
      指标一（单位） & 1.2 & 1.8 \\
      指标二（单位） & 2.4 & 3.0 \\
      \bottomrule
    \end{tabularx}
  \end{table}
  \source{表~\ref{tab:example} 使用占位数值。替换时可在此填写数据出处。}
\end{frame}

\begin{frame}{图形与图注}
  \begin{figure}
    \centering
    \begin{tikzpicture}[
      box/.style={draw=tsinghua, rounded corners=2pt, minimum width=2.3cm,
                  minimum height=.85cm, align=center},
      node distance=.85cm
    ]
      \node[box] (a) {环节一};
      \node[box, right=of a] (b) {环节二};
      \node[box, right=of b] (c) {环节三};
      \draw[tsinghua, thick, -{Stealth}] (a) -- (b);
      \draw[tsinghua, thick, -{Stealth}] (b) -- (c);
    \end{tikzpicture}
    % 使用自己的图时，可将 tikzpicture 换为下面这行，并填写实际文件名：
    % \includegraphics[width=.9\linewidth,height=.48\textheight,keepaspectratio]{assets/figure.pdf}
    \caption{流程示意}\label{fig:example}
  \end{figure}
  \source{图~\ref{fig:example} 为排版示意。替换时可在此填写图件出处。}
\end{frame}

% 不需要附录时，删除从 \appendix 到 \end{document} 之前的内容。
\appendix
\section{附录}

\begin{frame}{附录示例}
  在此填写补充材料。
\end{frame}

\end{document}
